% Generated by scripts/build_textbook.py; edit content/ and rebuild.

\chapter{Design and Defend an IT Service}

The final test of professional practice is not recalling a framework. It is putting a service design in front of people with different responsibilities, then answering their questions without hiding the costs, risks or unresolved choices. This chapter turns the preceding material into a practical design-and-defence exercise that can be completed alone or with a small team.



\subsection{The Service-Design Brief}


Design an IT service for a community organisation or small social-impact organisation and justify it end to end. A fictional organisation is acceptable when you need a safe portfolio exercise. If you work with a real organisation, obtain agreement from its operators, protect confidential information, and do not imply that consultation grants authority over community data. Work involving an Indigenous organisation or Indigenous data must be led by the relevant community authority; do not select such an organisation merely to make the project appear socially responsible.


A service is not an app. It is the complete arrangement that delivers continuing value: technology, support, incident handling, service promises, controlled change, lifecycle cost, vendor obligations and data authority. A proposal that presents software architecture and adds a support paragraph at the end has not designed a service.


The design pack should contain:


\begin{itemize}[leftmargin=*]
\item the organisation's purpose, users, constraints and decision rights;
\item a service model showing technology, people, suppliers and hand-offs;
\item intake, support, escalation, incident and change processes;
\item a delivery approach and a small set of measures that can drive improvement;
\index{licence}
\index{renewal}
\item vendor, licence and renewal assumptions;
\item a security and privacy baseline proportionate to the organisation;
\item a year-three operating cost, not merely an implementation price;
\item a data-authority position that states who decides about collection, access, reuse and disposal;
\item major risks, rejected alternatives and the evidence that would change the recommendation.
\end{itemize}

The frameworks in this book are lenses rather than mandatory ceremony. Explain why an ITIL-shaped support process earns its overhead for an organisation with one part-time technician, or why it does not. Explain what an error-budget discussion means when the on-call function is a volunteer. Show which delivery practices survive at this scale and which would be cargo cult. A strong design makes those choices visible.



\subsection{Build For A Mixed Review Panel}


Test the design against four perspectives: reliability and operations, service management, commercial sustainability, and community or data authority. Reviewers may be real colleagues or roles assigned during a rehearsal. The point is to make each perspective capable of blocking a weak design, not to stage a ceremonial approval.


The reliability reviewer asks what happens when the service fails outside office hours, who notices, and how it recovers. The service-management reviewer asks how a request becomes work, how work becomes controlled change, and where the authoritative record lives. The commercial reviewer asks who pays in year three, what happens at renewal, and how the organisation exits if a supplier raises prices or closes. The data-authority reviewer asks who decided the data could be collected, where it may go, who benefits, and what happens to it when the service winds down.


These perspectives model the review a service proposal encounters inside an organisation. A design that delights engineers can alarm a board; a design that reassures a board can leave operators with an impossible support promise. Prepare a ten-minute explanation that is credible to both.



\subsection{Run The Design Work In Deliberate Passes}


\index{backups}
Projects fail when contributors split the headings, write independently and combine the pieces at the end. Integration is a separate task. The cost model must match the architecture; the staffing assumption must match the service levels; the data-authority position must match the backup, analytics and supplier design. Schedule a whole-document review and trace every material promise to its owner, cost and failure mode.


Do not polish the presentation while the reasoning underneath remains weak. For every important claim, record the evidence, the strongest objection and what would change the decision. Run at least one full rehearsal with reviewers instructed to look for contradictions. Capture every question in a log, classify it as answered, design change or unresolved risk, and update the design rather than merely rehearsing a smoother answer.


Ground the work in operational constraints. For a real organisation, ask about the actual volunteer roster, funding cycle, support hours, tolerance for complexity and existing supplier contracts. For a fictional one, write those constraints before choosing technology. A constraint such as ``the treasurer will not store member data offshore'' changes the design; a generic promise to follow best practice does not.



\subsection{Defend The Design Under Questioning}


The question period rewards a specific temperament: candid, concrete and unhurried. When you do not know, say what is unknown and how you would resolve it. When a reviewer challenges an assumption, identify the trade-off and the evidence that would change your mind. That response demonstrates control of the design more convincingly than bluffing.


Keep answers short and name residual risk. A defensible proposal does not claim perfection. It shows that the designers understand the service's weaknesses, have assigned the important ones, and know which decisions require authority they do not possess.



\begin{quote}\small
The standard is not a perfect design. The standard is a design whose promises, costs, authority and failure modes can survive informed questions.
\end{quote}



\subsection{Turn The Result Into A Portfolio Artefact}


\index{DevOps}
\index{ITIL}
\index{vendor management}
After the defence, revise the pack once more. Add the question log, decision changes and unresolved risks. Remove confidential details and obtain permission before naming any real organisation or participant. The resulting artefact demonstrates more than familiarity with ITIL, DevOps or vendor management: it shows that you can reconcile operations, commercial constraints, technical delivery and data authority in one accountable service proposal.


The professional habit is portable. In later work you will sit in rooms where operators, salespeople, executives, suppliers and communities want different things. The purpose of this exercise is to arrive able to hold a position, expose its trade-offs and change it when the evidence demands.

\section{Practice Artefact}

Produce the service-design defence pack. It should contain the service model, support model, incident and change process, delivery approach, vendor and cost assumptions, data-authority position, year-three operating cost, and a question log from rehearsal.


The pack is ready when any team member can answer a panel question about a section they did not personally write.
